\chapter{まとめ}
\label{chap:conclusion}

本報告では，19\times19までのマインスイーパをIntel Cyclone IV E FPGA上で処理するシステムを示した．システムは，完全盤面を保持して開示を行う\texttt{minesweeper\_game\_core}，開示情報だけから推論する決定論ソルバ，メトリクス・表示回路，Virtual JTAG通信回路に分割されている．

アルゴリズム面では，数字セルの局所制約，盤面全体の残り地雷制約，部分集合差分，制約成分の限定全列挙，確率候補選択を組み合わせた．ハードウェア面では，FSMによる逐次化，同期RAMの待ち状態，固定容量FIFOとbitmap，乗算器を避けた19倍アドレス計算，交差積による確率比較，32クロックのrestoring divider，double-dabble表示変換を用いた．

通信面では，64ビットJTAGシフタ，トグル同期型CDCメールボックス，\texttt{jtag\_protocol\_engine}の\texttt{EX\_IDLE/BEGIN/DATA/COMMIT}，\texttt{board\_stream\_controller}の\texttt{ST\_IDLE/CFG/LOAD/COMMIT/RUN}を実装した．CRCと結果ACKを備えることで，盤面転送・計算・結果取得を分離し，通信エラーとソルバのstalledを区別できる．

実機のspeed10 1,000盤試験では，862盤を完全解答し，平均スコア0.9038487，平均81,110.615サイクルを得た．これはRTLをFPGAへ移植しただけでなく，入力プロトコル，逐次アルゴリズム，結果評価，表示，PC側検証を一つの再現可能な実験系としてまとめた成果である．

一方，最大12変数の列挙上限，固定容量の制約ストア，グローバル条件付けの未使用，確率選択のヒューリスティックは今後の課題である．これらを拡張すれば完全解答率を高められるが，RAM・論理資源・処理時間とのトレードオフが生じる．FPGA実装では，アルゴリズムの正しさだけでなく，観測可能情報，レジスタ幅，クロック当たりの処理量，通信の保持期間まで同時に設計する必要があることが，本成果物から得られる主要な知見である．
